\par\Needspace{4\baselineskip}
\SourceLeaf{020b}
\begin{center}
\pdfbookmark[2]{2.}{article-002}2.
\end{center}
Die eben dargelegten Gesetze unterscheiden sich von denen, die unser \textsc{Kepler} entdeckt hat, nur darin, dass sie in einer für alle Arten von Kegelschnitten geltenden Form angegeben sind und dass die Wirkung des bewegten Körpers auf die Sonne berücksichtigt wird, von der der Faktor $\BareRoot{(1+\mu)}$ abhängt. Betrachten wir diese Gesetze als Erscheinungen, die aus unzähligen und unzweifelhaften Beobachtungen abgeleitet sind, so zeigt uns die Geometrie, welche Wirkung
\SourceLeaf{021}die Sonne auf die sie umkreisenden Körper ausüben muss, damit diese Erscheinungen beständig auftreten. So findet man, dass die Sonne auf die sie umkreisenden Körper so wirkt, als triebe eine Anziehungskraft, deren Stärke dem Quadrat der Entfernung umgekehrt proportional ist, die Körper zum Mittelpunkt der Sonne hin. Gehen wir umgekehrt von der Annahme einer solchen Anziehungskraft als Prinzip aus, so ergeben sich diese Erscheinungen daraus als notwendige Folgen. Hier genügt es, die Gesetze angegeben zu haben. Bei ihrem Zusammenhang mit dem Gravitationsprinzip brauchen wir hier umso weniger zu verweilen, als nach dem großen \textsc{Newton} mehrere Autoren diesen Gegenstand behandelt haben, darunter der berühmte \textsc{Laplace}, der ihn in seinem vollendeten Werk \textit{Mécanique Céleste} so behandelt hat, dass nichts weiter zu wünschen blieb.

\begin{center}
\pdfbookmark[2]{3.}{article-003}3.
\end{center}

Untersuchungen über die Bewegungen der Himmelskörper, soweit sie in Kegelschnitten erfolgen, erfordern keineswegs eine vollständige Theorie dieser Kurvenart. Vielmehr genügt uns eine einzige allgemeine Gleichung, auf der alles Weitere aufbaut. Dabei scheint es am zweckmäßigsten, gerade diejenige Gleichung zu wählen, die sich bei der Untersuchung einer nach dem Anziehungsgesetz beschriebenen Kurve als charakteristische Gleichung ergibt. Bestimmen wir nämlich einen beliebigen Ort des Körpers auf seiner Bahn durch die Abstände $x, y$ von zwei in der Bahnebene gezogenen Geraden, die sich im Mittelpunkt der Sonne, also in einem der Brennpunkte der Kurve, rechtwinklig schneiden, und bezeichnen wir außerdem den stets positiv zu nehmenden Abstand des Körpers von der Sonne mit $r$, so erhalten wir zwischen $r, x, y$ die lineare Gleichung $r+\alpha x+\beta y=\gamma$. Darin bezeichnen $\alpha, \beta, \gamma$ konstante Größen, wobei $\gamma$ seiner Natur nach stets positiv ist. Ändern wir die an sich willkürliche Lage der Geraden, auf die sich die Abstände $x, y$ beziehen, so ändern sich die Form der Gleichung und der Wert von $\gamma$ offensichtlich nicht, sofern sich die Geraden weiterhin rechtwinklig schneiden. Dagegen nehmen $\alpha$ und $\beta$ andere Werte an. Offenbar kann man die Lage so bestimmen, dass $\beta$ den Wert $=0$ annimmt und $\alpha$ jedenfalls nichtnegativ ist. Schreiben wir dann für $\alpha, \gamma$ jeweils $e, p$, so nimmt unsere Gleichung die Form $r+ex=p$ an. Die Gerade, auf die sich dann die Abstände $y$ beziehen, heißt \emph{Apsidenlinie}, $p$ \emph{Halbparameter} und $e$ \emph{Exzentrizität}. Der Kegelschnitt schließlich heißt \emph{Ellipse}, \emph{Parabel} oder \emph{Hyperbel}, je nachdem, ob $e$ kleiner als, gleich oder größer als eins ist.

\SourceLeaf{022}Im Übrigen ist leicht einzusehen, dass die Lage der Apsidenlinie durch die angegebenen Bedingungen vollständig bestimmt ist, außer in dem einen Fall, in dem sowohl $\alpha$ als auch $\beta$ von vornherein $=0$ waren. In diesem Fall gilt stets $r=p$, auf welche Geraden sich die Abstände $x, y$ auch beziehen. Da also $e=0$ gilt, gehört die Kurve (die dann ein Kreis ist) nach unserer Definition zur Gattung der Ellipsen. Sie hat aber die Besonderheit, dass die Lage der Apsiden völlig beliebig bleibt, sofern man diesen Begriff auch auf diesen Fall ausdehnen will.
